\chapter{Optimal Stopping Problems}

Many decision problems under uncertainty can be formulated as the following question:
when should we stop observing a random process and take an action?
At each time $n$ we have collected information $\F_n$ (generated by the observations up to time $n$),
and we are allowed to stop only on the basis of that information (no ``peeking into the future'').
If stopping at time $n$ yields a random reward $Z_n$ (or incurs a cost), then an \emph{optimal stopping problem}
asks for a stopping time $T$ that maximizes the expected reward
\[
\sup_{T\le N}\E[Z_T]
\]
in a given time horizon $N$ (finite-horizon case).  The key subtlety is that the reward $Z_T$ depends on the
random time $T$, so one must optimize over a whole class of admissible random times.  Typical examples include:
selling an asset at the best time, deciding when to accept the first ``good'' candidate (secretary problem),
or stopping a sequential experiment once the evidence is strong enough.  The mathematical theory provides
two main tools: (i) the \emph{Snell envelope}, which formalizes the optimal ``stop vs.\ continue'' tradeoff via
backward recursion, and (ii) martingale/supermartingale methods, which yield sharp bounds and often identify
optimal stopping rules through hitting times of suitable regions.

\section{Basic definitions}
Let us first formalize some concepts needed for the discussion

\begin{definition}[Filtered probability space]
A \emph{filtered probability space} is a quadruple $(\Omega,\F,\Prob,(\F_n)_{0\le n\le N})$
where $N\in\N$ is fixed and
\[
\F_0 \subseteq \F_1 \subseteq \cdots \subseteq \F_N \subseteq \F
\]
is an increasing family of $\sigma$-algebras.
\end{definition}

\begin{definition}[Adapted process]
A stochastic process $Z=(Z_n)_{0\le n\le N}$ is called \emph{adapted} (to $(\F_n)$) if
$Z_n$ is $\F_n$-measurable for every $n$.
\end{definition}

\begin{definition}[Stopping time]
A random time $T:\Omega\to\{0,1,\dots,N\}$ is a \emph{stopping time} (w.r.t.\ $(\F_n)$) if
\[
\{T=n\}\in \F_n \qquad \text{for all } n=0,1,\dots,N.
\]
Equivalently, $\{T\le n\}\in\F_n$ for all $n$.
\end{definition}

\begin{remark}[Why stopping times?]
The condition $\{T=n\}\in\F_n$ formalizes the idea that at time $n$ we have enough
information to know whether we stop \emph{now}. A strategy that depended on future information
would not be implementable.
\end{remark}

When optimizing over (possibly uncountably many) stopping times, we need an intrinsically measurable notion
of the ``pointwise supremum'' of random variables.

\begin{definition}[Essential supremum]
Let $\{Z_t:t\in\mathcal{T}\}$ be a family of $\F$-measurable random variables (possibly indexed by an
uncountable set $\mathcal{T}$). A random variable $Z$ is the \emph{essential supremum} of the family,
written $Z=\esssup_{t\in\mathcal{T}} Z_t$, if
\begin{enumerate}[label=(\roman*)]
\item $Z$ is $\F$-measurable and $Z\ge Z_t$ a.s.\ for all $t\in\mathcal{T}$,
\item if $Z'$ is $\F$-measurable with $Z'\ge Z_t$ a.s.\ for all $t\in\mathcal{T}$, then $Z'\ge Z$ a.s.
\end{enumerate}
\end{definition}

\begin{remark}
The essential supremum is unique a.s.\ and always exists (one can construct it as the a.s.\ limit of
an increasing sequence of suprema over finite subfamilies).
\end{remark}

% =================================================
\section{Finite-horizon optimal stopping and the Snell envelope}

\subsection{Problem statement}

Let $Z=(Z_n)_{0\le n\le N}$ be an adapted process on $(\Omega,\F,\Prob,(\F_n))$ such that
$\E[|Z_n|]<\infty$ for all $n$. Think of $Z_n$ as the reward obtained by stopping at time $n$.

Let $\mathcal{T}_n$ denote the set of stopping times $T$ with values in $\{n,n+1,\dots,N\}$.
Define the \emph{value process} (or \emph{Snell envelope})
\begin{equation}\label{eq:Snell}
U_n := \esssup_{T\in\mathcal{T}_n}\E[Z_T\mid \F_n], \qquad 0\le n\le N.
\end{equation}
In particular, the optimal expected payoff at time $0$ is $u_0=\E[U_0]$.

\subsection{Backward recursion and optimal stopping time}

\begin{theorem}[Snell envelope recursion and existence of an optimal rule]\label{thm:snell}
Let $Z$ and $U$ be as above. Then:
\begin{enumerate}[label=(\roman*)]
\item $U_N=Z_N$ and for $n=1,\dots,N$,
\begin{equation}\label{eq:snell_recursion}
U_{n-1}=\max\bigl\{ Z_{n-1}, \E[U_n\mid \F_{n-1}] \bigr\}\quad\text{a.s.}
\end{equation}
\item For each $n$, the random time
\[
T_n^\ast := \min\{k\in\{n,\dots,N\}: U_k=Z_k\}
\]
is a stopping time in $\mathcal{T}_n$ and is optimal:
\begin{equation}\label{eq:optimality}
U_n=\E[Z_{T_n^\ast}\mid \F_n]\quad\text{a.s.}
\end{equation}
\end{enumerate}
\end{theorem}

\begin{remark}[Interpretation: stop vs.\ continue]
At time $n-1$ we can either:
\begin{itemize}
\item stop immediately and receive $Z_{n-1}$, or
\item continue; then the best we can do from time $n$ onward is $U_n$, hence your expected continuation
value is $\E[U_n\mid \F_{n-1}]$.
\end{itemize}
Therefore the optimal choice is the maximum of these two quantities, which is exactly
\eqref{eq:snell_recursion}.
\end{remark}

\begin{proof}
We prove both statements by backward induction.

\smallskip
\noindent\emph{Step 0: terminal time.}
At time $N$ we must stop, hence $\mathcal{T}_N=\{N\}$ and $U_N=\E[Z_N\mid\F_N]=Z_N$.

\smallskip
\noindent\emph{Step 1: recursion.}
Fix $n\in\{1,\dots,N\}$ and consider any $T\in\mathcal{T}_{n-1}$.
Split according to $\{T=n-1\}$ and $\{T\ge n\}$ (both events are in $\F_{n-1}$):
\begin{align*}
\E[Z_T\mid \F_{n-1}]
&= \E\bigl[\Ind_{\{T=n-1\}}Z_{n-1} + \Ind_{\{T\ge n\}} Z_T \,\bigm|\,\F_{n-1}\bigr] \\
&= \Ind_{\{T=n-1\}} Z_{n-1} + \Ind_{\{T\ge n\}} \E[Z_T\mid \F_{n-1}].
\end{align*}
Now use the tower property:
\[
\E[Z_T\mid \F_{n-1}]
= \Ind_{\{T=n-1\}} Z_{n-1} + \Ind_{\{T\ge n\}}\E\bigl[\E[Z_T\mid \F_n]\mid \F_{n-1}\bigr].
\]
Since on $\{T\ge n\}$ we have $T\in\mathcal{T}_n$, the definition of $U_n$ implies
$\E[Z_T\mid \F_n]\le U_n$ a.s., hence
\[
\E[Z_T\mid \F_{n-1}] \le \Ind_{\{T=n-1\}}Z_{n-1} + \Ind_{\{T\ge n\}}\E[U_n\mid\F_{n-1}]
\le \max\{Z_{n-1},\E[U_n\mid\F_{n-1}]\}.
\]
Taking the essential supremum over $T\in\mathcal{T}_{n-1}$ yields
$U_{n-1}\le \max\{Z_{n-1},\E[U_n\mid\F_{n-1}]\}$.
Conversely, both stopping immediately ($T\equiv n-1$) and continuing to an optimizer from time $n$
show that the right-hand side is attainable as an essential supremum, hence equality holds and we obtain
\eqref{eq:snell_recursion}.

\smallskip
\noindent\emph{Step 2: optimal stopping time.}
Define $T_n^\ast=\min\{k\ge n: U_k=Z_k\}$. Because $U$ and $Z$ are adapted,
\[
\{T_n^\ast = k\}=\Bigl(\bigcap_{j=n}^{k-1}\{U_j>Z_j\}\Bigr)\cap\{U_k=Z_k\}\in \F_k,
\]
so $T_n^\ast$ is a stopping time. One proves \eqref{eq:optimality} by induction using the recursion:
on $\{T_n^\ast=n\}$ we have $U_n=Z_n$; on $\{T_n^\ast\ge n+1\}$ we have $U_n=\E[U_{n+1}\mid \F_n]$
and then apply the induction hypothesis at time $n+1$.
\end{proof}

\begin{proposition}[Minimal supermartingale dominating $Z$]\label{prop:smallest_supermartingale}
The Snell envelope $U$ is a supermartingale and satisfies $U_n\ge Z_n$ for all $n$.
Moreover, if $M$ is any supermartingale with $M_n\ge Z_n$ for all $n$, then $M_n\ge U_n$ a.s.\ for all $n$.
\end{proposition}

\begin{proof}
From \eqref{eq:snell_recursion} we get $U_{n-1}\ge \E[U_n\mid\F_{n-1}]$, hence $U$ is a supermartingale.
Also $U_{n-1}\ge Z_{n-1}$ follows from the maximum. Minimality follows directly from the definition
\eqref{eq:Snell}: for any stopping time $T\in\mathcal{T}_n$ we have by optional sampling for supermartingales
$M_n\ge \E[M_T\mid\F_n]\ge \E[Z_T\mid\F_n]$, hence $M_n\ge \esssup_{T\in\mathcal{T}_n}\E[Z_T\mid\F_n]=U_n$.
\end{proof}

% =================================================
\section{Markovian optimal stopping and Bellman recursion}

The recursion in Theorem~\ref{thm:snell} is conceptually simple, but it involves conditional expectations
$\E[U_n\mid\F_{n-1}]$. In many applications one can reduce these to functions of a \emph{state variable}
using a Markov structure.

\subsection{Markov chain and transition kernel}

Let $(E,\mathcal{E})$ be a measurable space. A process $(X_n)_{0\le n\le N}$ is a \emph{time-homogeneous Markov chain}
with transition kernel $Q:E\times\mathcal{E}\to[0,1]$ if for all bounded measurable $\varphi:E\to\R$,
\[
\E[\varphi(X_{n+1})\mid \sigma(X_0,\dots,X_n)] = \E[\varphi(X_{n+1})\mid X_n]
= \int_E \varphi(y)\,Q(X_n,dy).
\]

\subsection{Bellman equation for Markovian rewards}

Assume now that the reward is of the form
\[
Z_n = g_n(X_n), \qquad 0\le n\le N,
\]
for measurable functions $g_n:E\to\R$ with $\E[|g_n(X_n)|]<\infty$.

\begin{theorem}[Dynamic programming / Bellman recursion]\label{thm:bellman}
Let $(X_n)$ be a Markov chain with kernel $Q$ and reward $Z_n=g_n(X_n)$. Then there exist measurable
value functions $V_n:E\to\R$ such that
\[
U_n = V_n(X_n)\quad\text{a.s.}
\]
and the functions satisfy the backward recursion
\begin{align}
V_N(x) &= g_N(x), \label{eq:bellman_terminal}\\
V_{n-1}(x) &= \max\Bigl\{ g_{n-1}(x), \int_E V_n(y)\,Q(x,dy)\Bigr\}, \qquad 1\le n\le N. \label{eq:bellman}
\end{align}
An optimal stopping time is given by
\[
T_n^\ast = \min\{k\ge n: g_k(X_k)=V_k(X_k)\},
\]
i.e.\ stop the first time we enter the \emph{stopping region}
$S_k:=\{x: g_k(x)=V_k(x)\}$.
\end{theorem}

\begin{proof}
We sketch the key idea. By induction on $n$:

For $n=N$ we have $U_N=Z_N=g_N(X_N)$ so $V_N=g_N$.
Assume $U_n=V_n(X_n)$ for some measurable $V_n$.
Then by Theorem~\ref{thm:snell},
\[
U_{n-1}=\max\{g_{n-1}(X_{n-1}),\,\E[V_n(X_n)\mid \F_{n-1}]\}.
\]
Since $X$ is Markov and $\F_{n-1}$ contains $\sigma(X_{n-1})$, the Markov property yields
\[
\E[V_n(X_n)\mid \F_{n-1}] = \E[V_n(X_n)\mid X_{n-1}]
= \int_E V_n(y)\,Q(X_{n-1},dy).
\]
Thus $U_{n-1}=V_{n-1}(X_{n-1})$ with $V_{n-1}$ defined by \eqref{eq:bellman}.
Optimality of the hitting time follows from Theorem~\ref{thm:snell}.
\end{proof}

% =================================================
\section{Martingale approach to finite-horizon optimal stopping}\label{sec:martingale_approach}

The Markovian approach in Section~6 reduces the recursion to functions of a state variable.
There is a complementary (and often very powerful) viewpoint based on \emph{martingales} and
\emph{supermartingale domination}.  It does not require a Markov structure and is therefore applicable
in very general filtered probability spaces.

Throughout this section we remain in discrete time with finite horizon $N$ and an integrable adapted
reward process $Z=(Z_n)_{0\le n\le N}$.

\subsection{The supermartingale domination principle}

\begin{definition}[Martingale / supermartingale]
An adapted integrable process $Y=(Y_n)_{0\le n\le N}$ is a \emph{martingale} if
$\E[Y_n\mid \F_{n-1}]=Y_{n-1}$ a.s.\ for all $n\ge 1$.
It is a \emph{supermartingale} if $\E[Y_n\mid \F_{n-1}]\le Y_{n-1}$ a.s.\ for all $n\ge 1$.
\end{definition}

\begin{proposition}[Domination gives an upper bound]\label{prop:domination_upper_bound}
Let $Y$ be a supermartingale such that $Y_n\ge Z_n$ a.s.\ for all $n$.
Then for every stopping time $T\le N$,
\[
\E[Z_T]\le \E[Y_T]\le \E[Y_0].
\]
Consequently,
\[
\sup_{T\le N}\E[Z_T] \le \E[Y_0].
\]
\end{proposition}

\begin{proof}
Since $Y_T\ge Z_T$, we have $\E[Z_T]\le \E[Y_T]$. Moreover, by optional sampling for supermartingales
in finite horizon, $\E[Y_T]\le \E[Y_0]$. Taking the supremum over $T$ yields the claim.
\end{proof}

\begin{remark}[How to use Proposition~\ref{prop:domination_upper_bound}]
If we can \emph{construct} a supermartingale $Y$ that dominates $Z$, we obtain a computable upper bound
$\E[Y_0]$ on the value. If we can also find a stopping time $T$ such that $Z_T=Y_T$ a.s.\ and
$\E[Y_T]=\E[Y_0]$ (i.e.\ equality in the optional sampling inequality), then $T$ must be optimal and
$\E[Y_0]=\sup_T\E[Z_T]$.
\end{remark}

\subsection{Snell envelope as the \texorpdfstring{smallest}{smallest} dominating supermartingale}

Section~5 already introduced the Snell envelope
\[
U_n=\esssup_{T\in\mathcal{T}_n}\E[Z_T\mid \F_n].
\]
From Proposition~\ref{prop:smallest_supermartingale} we know that $U$ is the \emph{minimal} supermartingale
dominating $Z$.  This provides an optimal stopping time and shows that the optimal value is attained.

\begin{theorem}[Martingale optimality principle]\label{thm:martingale_optimality_principle}
Let $U$ be the Snell envelope of $Z$ and let $T^\ast=\min\{n:U_n=Z_n\}$ (as in Theorem~\ref{thm:snell}).
Then:
\begin{enumerate}[label=(\roman*)]
\item $U$ is the smallest supermartingale such that $U_n\ge Z_n$ for all $n$.
\item For every stopping time $T\le N$, $\E[Z_T]\le \E[U_0]$.
\item The stopping time $T^\ast$ is optimal and $\E[Z_{T^\ast}]=\E[U_0]$.
\end{enumerate}
\end{theorem}

\begin{proof}
(i) is Proposition~\ref{prop:smallest_supermartingale}. (ii) follows from
Proposition~\ref{prop:domination_upper_bound} applied to $Y=U$.
For (iii), use Theorem~\ref{thm:snell} which shows $U_0=\E[Z_{T^\ast}]$.
\end{proof}

\subsection{Doob decomposition of the Snell envelope and the ``martingale part''}

A key martingale tool is the discrete-time Doob decomposition.

\begin{theorem}[Doob decomposition]\label{thm:doob_decomposition}
Let $Y=(Y_n)_{0\le n\le N}$ be an integrable adapted supermartingale. Then there exist:
\begin{itemize}
\item a martingale $M=(M_n)_{0\le n\le N}$ with $M_0=Y_0$,
\item a predictable, integrable, increasing process $A=(A_n)_{0\le n\le N}$ with $A_0=0$,
\end{itemize}
such that for all $n$,
\[
Y_n = M_n - A_n.
\]
Moreover, one can take increments
\[
\Delta A_n := A_n-A_{n-1} = Y_{n-1}-\E[Y_n\mid\F_{n-1}] \ge 0 \qquad (n\ge 1),
\]
and then define $M_n:=Y_n + A_n$.
\end{theorem}

\begin{proof}
Define $A_0=0$ and $\Delta A_n:=Y_{n-1}-\E[Y_n\mid\F_{n-1}]$ for $n\ge 1$, and set
$A_n=\sum_{k=1}^n \Delta A_k$. Then $A$ is increasing and predictable.
Let $M_n:=Y_n+A_n$. A direct computation shows
$\E[M_n\mid \F_{n-1}]=\E[Y_n\mid \F_{n-1}]+A_{n-1}+\Delta A_n=Y_{n-1}+A_{n-1}=M_{n-1}$,
so $M$ is a martingale and $Y=M-A$.
\end{proof}

Apply this to the Snell envelope $U$. Then $U=M-A$ with $M$ a martingale and $A$ increasing predictable.
The increments satisfy
\[
\Delta A_n = U_{n-1}-\E[U_n\mid\F_{n-1}]
= \max\{Z_{n-1},\E[U_n\mid\F_{n-1}]\} - \E[U_n\mid\F_{n-1}]
= (Z_{n-1}-\E[U_n\mid\F_{n-1}])^+.
\]
In particular, $\Delta A_n>0$ can only occur when it is optimal to stop at time $n-1$ (i.e.\ when
$U_{n-1}=Z_{n-1}$). This gives a precise meaning to the idea that the ``drift'' of the Snell envelope
is created only on the stopping region.

\subsection{A practical martingale recipe (finite horizon)}

The martingale approach is often used as follows.

\begin{enumerate}[label=\textbf{Step \arabic*.}]
\item \textbf{Guess a candidate value process.}
Propose an adapted process $Y$ that we believe equals the value $U$ (or at least dominates $Z$).
Typically one guesses $Y$ by problem structure (symmetry, convexity, candidate boundary, etc.).
\item \textbf{Verify supermartingale \& domination.}
Show that $Y_n\ge Z_n$ for all $n$ and $\E[Y_n\mid\F_{n-1}]\le Y_{n-1}$ for $n\ge 1$.
Then Proposition~\ref{prop:domination_upper_bound} yields $\sup_T\E[Z_T]\le \E[Y_0]$.
\item \textbf{Find a candidate optimal time.}
Define $T:=\min\{n: Y_n=Z_n\}$ (first hitting time of the contact set).
\item \textbf{Show tightness (attainment).}
Prove $\E[Y_T]=\E[Y_0]$ (often via a martingale argument on the stopped process) and $Y_T=Z_T$.
Then $\E[Z_T]=\E[Y_0]$, so $T$ is optimal and $Y=U$ (by minimality).
\end{enumerate}

% =================================================
\section{Example: coin tosses and the ``average number of heads'' game}

Consider i.i.d.\ Bernoulli random variables $(X_k)_{1\le k\le N}$ with
$\Prob(X_k=1)=\Prob(X_k=0)=\tfrac12$ and filtration $\F_n=\sigma(X_1,\dots,X_n)$ (and $\F_0=\{\varnothing,\Omega\}$).
Let $S_n=\sum_{k=1}^n X_k$ be the number of heads up to time $n$ and define the payoff
\[
Z_0=0,\qquad Z_n=\frac{S_n}{n}\quad(1\le n\le N).
\]
You may stop at any time $T\le N$ and receive $Z_T$.

\medskip
Because $S_n$ is a Markov chain on $\{0,1,\dots,n\}$, we can write $U_n=h_n(S_n)$ for suitable
functions $h_n$.

\begin{proposition}[Backward recursion for $h_n$]
There exist functions $h_n:\{0,1,\dots,n\}\to\R$ such that $U_n=h_n(S_n)$ and
\[
h_N(k)=\frac{k}{N},\qquad
h_{n-1}(k)=\max\Bigl\{\frac{k}{n-1},\,\frac12\bigl(h_n(k)+h_n(k+1)\bigr)\Bigr\}
\quad (2\le n\le N),
\]
with the convention $\frac{k}{0}=0$ for $n=1$.
\end{proposition}

\begin{proof}
Assume $U_n=h_n(S_n)$. Then
\[
\E[U_n\mid \F_{n-1}]
=\E[h_n(S_{n-1}+X_n)\mid \F_{n-1}]
=\frac12\bigl(h_n(S_{n-1})+h_n(S_{n-1}+1)\bigr),
\]
since $X_n$ is independent of $\F_{n-1}$. Insert this into \eqref{eq:snell_recursion} and define $h_{n-1}$
accordingly.
\end{proof}

\begin{remark}[Stopping boundary]
One always has $h_n(k)\ge k/n$ since $U_n\ge Z_n$. In fact $h_n$ is convex in $k$, so there is a threshold
$b_n^N$ such that stopping is optimal iff $S_n\ge b_n^N$:
\[
T_n^\ast=n \iff U_n=Z_n \iff h_n(S_n)=\frac{S_n}{n} \iff S_n\ge b_n^N.
\]
\end{remark}

% =================================================
\section{Example: the secretary problem}

There are $N$ candidates arriving sequentially. Exactly one candidate is the best overall, but we can only
observe \emph{relative ranks} among the candidates seen so far. Once we reject a candidate, we cannot recall them.

Let $(Y_1,\dots,Y_N)$ be a random permutation of $\{1,\dots,N\}$, where $Y_n=1$ means ``candidate $n$ is best overall''.
At time $n$ we observe only whether candidate $n$ is best among the first $n$.
Define the \emph{relative rank}
\[
X_n := \#\{1\le k\le n : Y_k \le Y_n\}\in\{1,\dots,n\}.
\]
Then $X_n=1$ means ``candidate $n$ is the best so far''. One can show that $X_n$ is uniform on $\{1,\dots,n\}$, and
that $(X_n)$ is a Markov chain (in fact, it is the classical \emph{relative rank process}).

\medskip
We wish to maximize the probability of selecting the best overall:
\[
\sup_{T\le N} \Prob(Y_T=1),
\]
where $T$ ranges over stopping times w.r.t.\ the information $\F_n^X:=\sigma(X_1,\dots,X_n)$.

\begin{proposition}[Payoff representation]
For any stopping time $T$ w.r.t.\ $(\F_n^X)$,
\[
\Prob(Y_T=1) = \E\bigl[f_T(X_T)\bigr],
\qquad
f_n(x):=\begin{cases}
\frac{n}{N}, & x=1,\\
0, & x\ne 1.
\end{cases}
\]
\end{proposition}

\begin{proof}
Condition on $\F_n^X$. If $X_n\ne 1$, the $n$-th candidate is not best so far, hence cannot be best overall,
so $\Prob(Y_n=1\mid \F_n^X)=0$. If $X_n=1$, then among the first $n$ candidates, candidate $n$ is best.
By symmetry of the random permutation, the event ``the best overall is among the first $n$'' has probability $n/N$,
and on $\{X_n=1\}$ this is exactly the same as $\{Y_n=1\}$. Thus $\Prob(Y_n=1\mid \F_n^X)=\frac{n}{N}\Ind_{\{X_n=1\}}$,
and taking expectation at $n=T$ yields the claim.
\end{proof}

Therefore the secretary problem becomes a finite-horizon optimal stopping problem for the adapted reward
process $Z_n:=f_n(X_n)$ because the \emph{only} decision we make is a stopping time $T\in\{1,\dots,N\}$
with respect to the available information $\F_n^X=\sigma(X_1,\dots,X_n)$, and the performance criterion
\[
\Prob(Y_T=1)
\]
can be written solely in terms of that stopping time and the observed process $(X_n)$:
by the previous proposition,
\[
\Prob(Y_T=1)=\E\!\left[f_T(X_T)\right]=\E[Z_T].
\]
Hence maximizing the success probability over all admissible decision rules is \emph{equivalent} to
maximizing $\E[Z_T]$ over all $\F^X$-stopping times $T\le N$, which is exactly the standard finite-horizon
optimal stopping problem for the adapted payoff process $(Z_n)_{n=1}^N$.

\subsection*{Snell envelope and the threshold rule}

\paragraph{Plan of attack.}
After Proposition~9.1 we have reduced the original goal
\[
\sup_{T\le N}\Prob(Y_T=1)
\]
to a genuine optimal stopping problem for the adapted payoff process $Z_n=f_n(X_n)$, i.e.
\[
\sup_{T\le N}\E[Z_T].
\]
Our goal now is to \emph{compute}
(or at least characterize) an optimal stopping time. The strategy is:
\begin{itemize}
\item First we introduce the \emph{Snell envelope} $U_n$, which represents the best achievable
\emph{conditional} success probability when we stand at time $n$ and have observed the history.
By definition, $U_n$ compares all future stopping rules and hence encodes the optimal continuation value.
\item Then we use the Snell-envelope recursion
\[
U_n=\max\{Z_n,\E[U_{n+1}\mid \F_n^X]\},
\]
which expresses the optimal decision at time $n$ as a simple \emph{stop vs.\ continue} comparison:
stop now and get $Z_n$, or continue and get the conditional expected value of optimally stopping later.
\item Finally, we exploit the special structure of the secretary problem:
the payoff $Z_n$ is positive only on the event $\{X_n=1\}$ (``best so far'').
This collapses the recursion to a one-dimensional comparison between a known immediate payoff $n/N$
and a time-dependent continuation value.  We show that this comparison switches only once,
which leads to a \emph{threshold rule}: ignore the first $m-1$ candidates, then accept the first
candidate who is best so far.
\end{itemize}
In short: we compute optimal stopping by repeatedly comparing \emph{immediate reward} and
\emph{expected continuation value}, and the structure of the secretary problem forces the optimal rule
to be a simple time threshold.


Let
\[
U_n := \esssup_{T\in\mathcal{T}_n}\E[f_T(X_T)\mid \F_n^X].
\]
Because of the Markov structure, $U_n$ depends on the past only through $X_n$, so we write
$U_n = u_n(X_n)$.

A key simplification is that it is never optimal to stop when $X_n\ne 1$, since then the payoff is $0$
and we can always do at least as well by continuing. Hence only the state $x=1$ matters.

\begin{lemma}[Recursion in terms of continuation values]\label{lem:sec_rec}
Let $c_n:=\E[U_{n+1}\mid \F_n^X]$ be the continuation value at time $n$. Then
\[
U_n = \max\{f_n(X_n),c_n\}.
\]
Moreover, using the fact that $X_{n+1}$ is uniform on $\{1,\dots,n+1\}$ conditionally on $\F_n^X$, one obtains
\[
c_n = \frac{1}{n+1}\sum_{i=1}^{n+1} u_{n+1}(i).
\]
\end{lemma}

\begin{proof}
Recall that $Z_n=f_n(X_n)$ and that the Snell envelope (with respect to the filtration
$\F_n^X=\sigma(X_1,\dots,X_n)$) is defined by
\[
U_n:=\esssup_{T\in\mathcal{T}_n}\E[Z_T\mid \F_n^X],
\]
where $\mathcal{T}_n$ denotes the set of $\F^X$-stopping times taking values in $\{n,n+1,\dots,N\}$.
By the general Snell-envelope recursion (Theorem~\ref{thm:snell}), we have for $n=0,1,\dots,N-1$,
\begin{equation}\label{eq:snell_sec}
U_n=\max\{Z_n,\E[U_{n+1}\mid \F_n^X]\}
=\max\{f_n(X_n),c_n\},
\end{equation}
where we set $c_n:=\E[U_{n+1}\mid \F_n^X]$. This proves the first claim.

\medskip
For the second claim, we use the Markov/symmetry structure of the relative-rank process.
Since $U_{n+1}$ is $\F_{n+1}^X$-measurable and (by the Markovian reduction) can be written as
\[
U_{n+1}=u_{n+1}(X_{n+1})
\]
for some function $u_{n+1}:\{1,\dots,n+1\}\to\R$, it follows that
\begin{equation}\label{eq:cn_def}
c_n=\E[U_{n+1}\mid \F_n^X]=\E[u_{n+1}(X_{n+1})\mid \F_n^X].
\end{equation}
Now we claim that, conditionally on $\F_n^X$, the next relative rank $X_{n+1}$ is uniformly distributed
on $\{1,\dots,n+1\}$. Indeed, given the relative order information up to time $n$, the position of the
$(n+1)$-st observation among the first $n+1$ observations is equally likely to be any of the ranks
$1,2,\dots,n+1$ (this is a standard symmetry property of random permutations / i.i.d.\ continuous samples).
Therefore,
\[
\Prob(X_{n+1}=i\mid \F_n^X)=\frac{1}{n+1},\qquad i=1,\dots,n+1,
\]
and inserting this into \eqref{eq:cn_def} yields
\[
c_n
=\sum_{i=1}^{n+1} u_{n+1}(i)\,\Prob(X_{n+1}=i\mid \F_n^X)
=\frac{1}{n+1}\sum_{i=1}^{n+1}u_{n+1}(i),
\]
which is exactly the stated formula.
\end{proof}


\begin{remark}
The last identity is a concrete instance of ``reducing $\F_n$ to crucial information'':
the conditional distribution of $X_{n+1}$ given the past does not depend on the full past.
\end{remark}

\paragraph{Idea of Lemma~\ref{lem:sec_rec}.}
Lemma~\ref{lem:sec_rec} isolates the \emph{only} quantity we really need in order to apply the Snell-envelope recursion
to the secretary problem: the \emph{continuation value} $c_n=\E[U_{n+1}\mid\F_n^X]$, i.e.\ the optimal
success probability \emph{if we decide not to stop at time $n$}.  The first statement,
$U_n=\max\{f_n(X_n),c_n\}$, is just the general ``stop vs.\ continue'' principle: at time $n$ we either
stop immediately and obtain the payoff $f_n(X_n)$ (which is $n/N$ if the current candidate is best so far,
and $0$ otherwise), or we continue and then your best possible performance from the next step onward is
$U_{n+1}$, whose value today is its conditional expectation $c_n$.  The second statement explains how to
\emph{compute} this continuation value: conditionally on the information up to time $n$, the next relative
rank $X_{n+1}$ is uniformly distributed on $\{1,\dots,n+1\}$ (by symmetry of the random permutation), so
the conditional expectation of $U_{n+1}=u_{n+1}(X_{n+1})$ becomes the simple average
$c_n=\frac{1}{n+1}\sum_{i=1}^{n+1}u_{n+1}(i)$.  Thus Lemma~9.2 turns the abstract conditional expectation
$\E[U_{n+1}\mid\F_n^X]$ into an explicit quantity and thereby makes the backward induction feasible.


A classical analysis shows that the optimal strategy is a \emph{threshold rule}:
ignore a certain number of first candidates, then pick the first candidate who is best so far.

\begin{theorem}[Threshold rule and asymptotics]
There exists an index $m\in\{1,\dots,N\}$ such that an optimal strategy is:
\begin{quote}
Reject candidates $1,\dots,m-1$, then stop at the first time $n\ge m$ with $X_n=1$.
\end{quote}
Moreover, $m$ is characterized by the harmonic sum condition
\[
\sum_{k=m}^{N-1}\frac{1}{k}\le 1 < \sum_{k=m-1}^{N-1}\frac{1}{k},
\]
and therefore $m\sim \frac{N}{e}$ as $N\to\infty$. In particular, the optimal success probability converges to $1/e$.
\end{theorem}

\begin{proof}[Idea of proof]
The recursion from Lemma~\ref{lem:sec_rec} implies that stopping at time $n$ with $X_n=1$ is optimal iff
\[
\frac{n}{N} \ge c_n.
\]
One proves by backward induction that the continuation values can be written in terms of harmonic sums,
and that the inequality above switches from false to true at a single index $m$.
Finally, compare harmonic sums with logarithms:
\[
\sum_{k=m}^{N-1}\frac{1}{k}\approx \int_m^N \frac{dx}{x}=\log\!\Bigl(\frac{N}{m}\Bigr),
\]
so the threshold condition is approximately $\log(N/m)\approx 1$, i.e.\ $m\approx N/e$.
\end{proof}


